% Standalone problem document, generated from the adjacent README.md.
% Compile with XeLaTeX. Mathematical content is maintained in Markdown.
\documentclass[11pt,a4paper]{article}
\usepackage[fontsize=10.5pt]{scrextend}
\usepackage[margin=22mm,headheight=15pt,headsep=9mm,footskip=11mm]{geometry}
\usepackage{fontspec}
\setmainfont{texgyrepagella}[Extension=.otf,UprightFont=*-regular,BoldFont=*-bold,ItalicFont=*-italic,BoldItalicFont=*-bolditalic]
\setsansfont{texgyreheros}[Extension=.otf,UprightFont=*-regular,BoldFont=*-bold,ItalicFont=*-italic,BoldItalicFont=*-bolditalic]
\setmonofont{texgyrecursor}[Extension=.otf,UprightFont=*-regular,BoldFont=*-bold,ItalicFont=*-italic,BoldItalicFont=*-bolditalic,Scale=MatchLowercase]
\usepackage{amsmath,amssymb}
\usepackage{unicode-math}
\setmathfont{texgyrepagella-math.otf}
\usepackage{xcolor,microtype,enumitem,needspace,fancyhdr}
\usepackage{bookmark}
\definecolor{ink}{HTML}{17344B}
\definecolor{accent}{HTML}{197D80}
\definecolor{muted}{HTML}{596773}
\hypersetup{colorlinks=true,linkcolor=accent,urlcolor=accent,pdftitle={IE-07 - Deterministic
regularization of the nonsymmetric
eigenproblem},pdfauthor={Open Problems in Numerical Linear Algebra}}
\urlstyle{same}
\setlength{\parindent}{0pt}
\setlength{\parskip}{5pt plus 1pt minus 1pt}
\linespread{1.04}
\setlength{\emergencystretch}{3em}
\widowpenalty=10000
\clubpenalty=10000
\displaywidowpenalty=10000
\allowdisplaybreaks[1]
\setlist{nosep,leftmargin=1.5em}
\providecommand{\tightlist}{\setlength{\itemsep}{0pt}\setlength{\parskip}{0pt}}
\providecommand{\pandocbounded}[1]{#1}
\pagestyle{fancy}
\fancyhf{}
\fancyhead[L]{\sffamily\footnotesize\color{muted}OPEN PROBLEMS IN NUMERICAL LINEAR ALGEBRA}
\fancyhead[R]{\sffamily\bfseries\footnotesize\color{accent}IE-07}
\fancyfoot[L]{\parbox[b]{0.9\textwidth}{\sffamily\fontsize{8}{10}\selectfont\color{muted}Editorial ratings; a bounded search cannot certify that a problem remains unsolved.\\Literature check: 2026-09-10}}
\fancyfoot[R]{\sffamily\footnotesize\color{muted}\thepage}
\renewcommand{\headrulewidth}{0.3pt}
\renewcommand{\headrule}{\hbox to\headwidth{\color{accent}\leaders\hrule height \headrulewidth\hfill}}
\makeatletter
\renewcommand{\section}{\Needspace{5\baselineskip}\@startsection{section}{1}{0pt}{12pt plus 2pt minus 2pt}{4pt}{\sffamily\bfseries\large\color{ink}}}
\renewcommand{\subsection}{\Needspace{4\baselineskip}\@startsection{subsection}{2}{0pt}{9pt plus 1pt minus 1pt}{3pt}{\sffamily\bfseries\normalsize\color{ink}}}
\makeatother
\setcounter{secnumdepth}{0}
\begin{document}
{\sffamily\small\color{muted}Eigenvalues and inverse problems\par}
\vspace{3pt}
{\raggedright\hyphenpenalty=10000\exhyphenpenalty=10000\sffamily\bfseries\fontsize{21}{24}\selectfont\color{ink}Deterministic
regularization of the nonsymmetric eigenproblem\par}
\vspace{7pt}
{\sffamily\small\textbf{Difficulty:} extreme\quad\textbf{Importance:} broadly
interesting\par}
{\sffamily\small\bfseries\color{accent}Status: Open\par}
\vspace{4pt}
{\color{accent}\hrule height 0.8pt}
\vspace{6pt}
\textbf{Rating rationale:} Extreme reflects deterministic regularization
of a fully general nonnormal eigenproblem; broad impact comes from
replacing randomness in a basic numerical linear algebra primitive.

\subsection{Problem statement}\label{problem-statement}

For a matrix \(M\) with simple spectrum, define

\[
\operatorname{gap}(M)=\min_{i\ne j}|\lambda_i(M)-\lambda_j(M)|,
\qquad
\kappa_V(M)=\inf_{M=VDV^{-1},\ D\text{ diagonal}}\|V\|_2\|V^{-1}\|_2.
\]

For every \(n\geq2\), \(A\in\mathbb C^{n\times n}\) with
\(\|A\|_2\leq1\), and \(0<\delta<1/2\), construct deterministically a
matrix \(E\) such that

\[
\|E\|_2\leq\delta,\qquad A+E\text{ has simple spectrum},\qquad
\frac{\kappa_V(A+E)}{\operatorname{gap}(A+E)}\leq C(n/\delta)^c,
\]

using \(O(n^3\log^d(n/\delta))\) exact arithmetic operations, for
universal constants \(C,c,d\). The task is to find the perturbation, not
merely prove its existence.

\subsection{References}\label{references}

Banks, Garza-Vargas, Kulkarni, and Srivastava,
\href{https://doi.org/10.1007/s10208-022-09577-5}{\emph{Pseudospectral
Shattering, the Sign Function, and Diagonalization in Nearly Matrix
Multiplication Time}}, FOCM 23 (2023), §6, first future-research
question. Amsel et al.,
\href{https://arxiv.org/html/2602.05394v3\#S3.SS1}{\emph{Linear Systems
and Eigenvalue Problems}}, August 2026, Problem 3.1.

\subsection{Earlier status check ---
2026-09-08}\label{earlier-status-check-2026-09-08}

Searches for \textit{deterministic pseudospectral shattering 2026}
found randomized and exponential-bound deterministic results, not the
stated algorithm. The normalization and exponent orientation here are
explicit: Problem 3.1's printed \((\delta/n)^c\) contradicts its own
preceding motivation; \((n/\delta)^c\) is the intended polynomial upper
bound.

\subsection{Audit update --- 2026-09-10}\label{audit-update-2026-09-10}

Rechecked \href{https://arxiv.org/html/2602.05394v3}{workshop version
3}, Problem 3.1. Deterministic pseudospectral-regularization searches
found no solution. The corrected reciprocal scale in this entry remains
necessary; the source's displayed exponent does not match the intended
inverse-polynomial regularity guarantee.
\end{document}
